\documentclass[11pt,a4paper]{article}
\usepackage[utf8]{inputenc}
\usepackage[T1]{fontenc}
\usepackage[spanish,es-noshorthands]{babel}
\usepackage{lmodern}
\usepackage{amsmath,amssymb}
\usepackage[margin=2.5cm]{geometry}
\usepackage{microtype}
\usepackage{enumitem}
\usepackage[hidelinks]{hyperref}
\setlength{\parskip}{0.5em}
\setlength{\parindent}{0pt}
\setlist{noitemsep,topsep=2pt}

\newcommand{\CBm}{$\mathrm{CBSOM}^{-}$}
\newcommand{\CBp}{$\mathrm{CBSOM}^{+}$}
\newcommand{\SOMm}{$\mathrm{SOM}^{-}$}
\newcommand{\pag}[1]{\textit{(pág.~#1)}}


\begin{document}

La nueva versión del paper sigue la estructura del artículo de Fairness
(\textit{Fairness and Consensus in an Asynchronous Opinion Model for Social
Networks}). Las páginas corresponden al PDF nuevo, que tiene 21; el anterior
tenía 20.

\section*{Abstract \pag{1}}

El abstract sigue el mismo orden del de Fairness: qué modelo presentamos, cómo
funciona un paso, qué probamos y qué muestran las simulaciones. Cuenta en
palabras lo que hace cada agente: pasa su desacuerdo con cada influenciador por
la función de sesgo, actualiza su opinión y después decide si habla o se queda
callado. También deja claro que en \CBm{} los callados no influyen a nadie,
mientras que en \CBp{} siguen influyendo con la última opinión que dijeron en
público.

Salieron los detalles de implementación del simulador y la frase final sobre
estrategias contra la polarización. Para backfire y authority, el abstract dice
ahora que impiden el acuerdo ``almost always'' en vez de ``robustly'', porque
authority tiene eventos sueltos con $n=10$. Quedó en 197 palabras; CLEI pide
máximo 200.

\section*{Introducción \pag{1--2}}

El primer párrafo arranca como el de Fairness, con el peso de las redes
sociales en la formación de opiniones, y trae oraciones más cortas. En el de
DeGroot, ``repetitively'' pasó a ``repeatedly'', y la frase sobre consenso quedó
como en Fairness: el consenso es un problema central del aprendizaje social, y
no alcanzarlo indica una comunidad polarizada. Los párrafos de sesgos y de
Alvim et al.\ ahora son uno solo, sin la frase en negrita sobre la ``central
contribution''.

Entran tres párrafos nuevos. ``In this paper, we combine\ldots'' explica el
modelo sin fórmulas. El estado guarda la opinión de cada agente y si está
hablando o callado; en cada paso el agente filtra cada desacuerdo con su sesgo,
revisa su opinión y decide si habla en la siguiente ronda según su radio de
tolerancia. Ese párrafo aclara además que los callados siguen actualizando su
opinión, solo que no la expresan. El segundo cuenta los resultados en orden:
propiedades estructurales bajo sesgos en $R$, consenso en cliques para
$f_\alpha$ (con $n=2$ incluido) y lo que exploran las simulaciones por fuera de
la teoría. El tercero adelanta el ``To the best of our knowledge'', que antes
aparecía solo en las conclusiones.

En las contribuciones, el ítem de simulaciones hablaba de una interacción entre
tamaño, topología y sesgo ``in nuanced ways''. Ahora nombra los hallazgos de la
Discusión: el umbral de densidad crece de forma sub-lineal con el tamaño de la
red y la memoria casi acaba con el consenso en redes grandes. Organization
cierra con dos datos nuevos: varias pruebas aparecen esbozadas en el texto y
completas en el Apéndice~A, y el simulador tiene su link a GitHub.

\section*{Modelos \pag{2--6}}

Cada definición trae ahora un par de líneas que dicen qué significa. Después de
la Definición~1 aparece qué son los vértices, qué quiere decir una arista
$(j,i)$, qué mide el peso y qué exige la Ecuación~(1). En la parte de opiniones
está lo que significan $B_i^t=0$ y $B_i^t=1$. Tras la Definición~2,
$\beta_{i,j}(x)$ queda explicada como la parte del desacuerdo que el agente
absorbe de su influenciador. La presentación de la región $R$ dice por qué
importa: un sesgo en $R$ acerca al agente a su influenciador sin pasarse.

La frase ``Throughout this paper all bias functions lie in $R$'' chocaba con los
experimentos, que usan backfire y authority. Quedó como ``Our formal results
assume bias functions in $R$''. En \SOMm{} hay una línea nueva: los agentes
callados siguen actualizando su opinión.

El cambio más visible es el Example~1 \pag{5}, justo después de la definición
de \CBm. Es un clique de 3 agentes con pesos $1/2$,
$\mathbf{B}^0=(1,0.5,0)$, $\tau_i=0.4$ y el sesgo $\beta(x)=x/2$. Las
ecuaciones del modelo dan $\mathbf{B}^1=(0.625,0.5,0.375)$; sin sesgo, los
agentes 1 y 3 se cruzan y terminan en $(0.25,0.5,0.75)$. Después todos se
callan, las opiniones se quedan quietas una ronda, todos vuelven a hablar y
convergen a $0.5$. Con eso el lector ve en un caso concreto lo que prueba el
Lema~5: el silencio global dura máximo una ronda. Las cuentas están verificadas
con un script, y \texttt{main.tex} tiene un \verb|\newtheorem{example}| para el
ejemplo.

Las oraciones que empezaban con el nombre del modelo arrancan ahora con ``The
model\ldots'', y ``For any'' pasó a ``For every''.

\section*{Análisis \pag{6--10}}

La sección abre con los dos teoremas principales y la ruta de la prueba en tres
pasos, como en el de Fairness. Primero, silencio y sesgo no se estorban y los
extremos convergen a $U$ y $L$. Luego, $f_\alpha$ contrae la dispersión en
cliques. Al final, una contradicción da $U=L$. El caso $n=2$ va aparte. Entre
el Corolario~6 y el Lema~7 hay una frase con la intuición de ambos.

La notación cambió en dos puntos. La dispersión se llama ahora $\Delta_t$: $R$
ya nombraba la región receptive-resistant, y una frase como ``since
$\beta\in R$ \ldots\ $R_t<\varepsilon$'' se leía rara. El cambio llega también
al apéndice. El cuadrado $[-1,1]^2$ dejó de llamarse $S$, letra que ya usa el
vector de silencio.

El enunciado del Lema~11 definía $\alpha_{\min}$ e $I_{\min}$ sin usarlos;
ahora esas definiciones aparecen en el esbozo de la prueba, donde sí hacen
falta. La leyenda de la Figura~3 compara consenso con consenso: ``consensus is
reached faster than under $\mathrm{conf}(x)$''. En el caso diádico, ``\SOMm{}
fails to converge'' quedó como ``can fail to converge'', porque la oscilación
depende de los parámetros.

El Remark~3.3 decía que un $\alpha$ grande oscila y uno pequeño no. Ahora da la
condición exacta: con $2\alpha I>1$, $\gamma<0$ y la dinámica oscila mientras
converge; con $2\alpha I<1$, la convergencia es monótona. Eso cuadra con las
Figuras~5 y~6 y con el caso $\alpha=0.5$. Además desapareció ``Case~2 of
Theorem~14'' del texto y de la leyenda de la Figura~5: ese teorema no tiene
casos.

\section*{Experimentos \pag{10--16}}

La introducción de la sección presenta las simulaciones como el estudio de lo
que la teoría no alcanza a cubrir: grafos libres de escala, sesgos fuera de
$R$, mezclas y memoria. La versión anterior decía que ``corroboraban'' los
teoremas.

El resto es sobre todo redacción. Las oraciones largas quedaron partidas, en
especial la del umbral sub-lineal, que pasaba de 50 palabras. El párrafo sobre
por qué escogimos esas proporciones ahora son tres: por qué son heurísticas,
qué pregunta responde cada configuración y qué queda para trabajo futuro. Los
encabezados de párrafo tienen todos el mismo formato.

La configuración~(i) \pag{13} tiene un ajuste de contenido. El texto la
describía como ``an $R$-region majority'' con ``a 70\% $R$-region fraction'',
pero la lista de la misma sección la define como 70\% sin sesgo y 30\% de
confirmación. Ahora dice que mezclar aristas con confirmación y sin sesgo en
proporción 30/70 conserva el consenso.

El tono bajó en varios puntos. Salieron ``Remarkably, a mere\ldots'',
``consensus killers'' y ``sharp''. La frase sobre la ``structural
impossibility'' que causaba la memoria dice ahora que a gran escala no
observamos consenso, que es lo que alcanza a mostrar una simulación. En la
Discusión, ``Confirmation reduces the threshold by one unit compared to
\SOMm'' quedó como ``the threshold under Confirmation is one unit lower than
under \SOMm''.

En densidad relativa, los párrafos se llamaban ``Row~1'' a ``Row~4'', aunque
cada figura tiene dos filas. Ahora cada uno dice a qué figura y fila se
refiere, por ejemplo ``Figure~11, top row''.

\section*{Conclusiones \pag{16--17}}

Las dos subsecciones se volvieron una sola sección, con el orden del de
Fairness: resumen, trabajo relacionado y trabajo futuro. El resumen aclara que
las propiedades estructurales están probadas para el modelo sin memoria.

El párrafo de Alvim et al.\ arranca con ``The work closest to ours'' y dice en
qué nos diferenciamos: nuestro resultado cubre solo cliques y $f_\alpha$. Salió la frase que
planteaba como problema abierto extender ese resultado a grafos arbitrarios. En el
de Aranda et al., ``confirms that memory acts as a structural barrier'' pasó a
``suggests'', porque esa evidencia sale de simulaciones.

Hay un párrafo nuevo que compara el radio de tolerancia con los modelos de
confianza acotada de Deffuant et al. Allá el umbral decide a quién escucha el
agente; acá decide si habla. El mismo párrafo ubica el trabajo frente a las
simulaciones de espiral del silencio que ya citábamos (Sohn, Wang, Ross,
Cabrera). Deffuant estaba en el \texttt{.bib} sin citar y ahora aparece en las
referencias.

El trabajo futuro quedó en dos puntos: llevar los sesgos cognitivos a los
modelos de silencio basados en confianza de Gaona, que ahora se cita con su
tesis de maestría, y estudiar topologías de red dinámicas. Salieron la
extensión a grafos arbitrarios y la validación con datos empíricos.

\section*{Apéndice \pag{18--21}}

El apéndice usa $\Delta_t$ en lugar de $R_t$. También cambió ``analyse'' por
``analyze'', para mantener ortografía estadounidense en todo el paper, y ``Fix
any agent'' por ``Fix an agent''.

\end{document}
